\documentclass[11pt]{article}
\title{Testing Stage-Dependent Injury-to-Regeneration Coupling in Xenopus laevis: A Hypothesis-Led Study Plan}
\date{}
\author{SIGMA AI-assisted research synthesis}
\usepackage[T1]{fontenc}
\usepackage[utf8]{inputenc}
\usepackage[margin=1in]{geometry}
\usepackage{amsmath,amssymb,booktabs,longtable}
\usepackage{microtype}
\usepackage{titlesec}
\emergencystretch=2em
\raggedbottom
\linespread{1.07}
\usepackage{xcolor}
\definecolor{blue}{HTML}{24485E}
\titleformat{\section}{\Large\bfseries\sffamily\color{blue}}{\thesection}{0.8em}{}
\titleformat{\subsection}{\large\bfseries\sffamily}{\thesubsection}{0.7em}{}
\setlength{\parskip}{.6em}
\setlength{\parindent}{0pt}
\begin{document}
\maketitle

\begin{abstract}
The developmental loss and subsequent recovery of tail-regeneration competence in Xenopus laevis raises a mechanistic question: does refractory tissue fail to enter a regenerative epidermal state, or fail to translate an early injury response into that state? This article would prioritize a stage-dependent coupling hypothesis rather than treat regeneration-initiating cells (RICs), extracellular-matrix (ECM) remodeling, or familiar injury pathways as new discoveries. The candidate prediction is that early wound-response intensity has a weaker positive relationship with subsequent epidermal RIC/ECM activation at stage 46 than at stage 40.

The proposed analysis would first establish whether the public transcriptomic samples permit that relationship to be estimated, then compare it with stage-only, RIC/ECM-only, and known-signaling baselines while addressing biological replication and tissue composition. The supplied evidence contains no completed analysis, numerical effect estimate, independently validated predictor, or perturbation result. Accordingly, the planned article is a falsifiable hypothesis and analysis framework, not a report of a newly discovered mechanism; novelty, animal-level prediction, and the explanation of later recovery remain unresolved.
\end{abstract}
\section*{Evidence coverage limitation}
The evidence index omitted 1863 material saved research field fragment(s) across 6 source(s). These fragments were not reviewed for this paper. Interpretations that could depend on omitted material remain provisional; the complete frozen source snapshot and omission ledger are retained with the generation record.
\section{Introduction and biological context}
The target phenomenon is a developmental change in tail-regeneration competence within Xenopus laevis, including variation between animals and the later return of competence. A comparison between stage classes can help locate a developmental difference, but it cannot by itself explain individual regeneration outcomes or the subsequent recovery of competence.

The project brief identifies transient basal-epidermis-derived RICs and defective ECM remodeling as prior findings, alongside established interest in injury, inflammatory, developmental, and metabolic signals. These observations define the background against which an interaction must be tested; they are not the proposed contribution. The supplied material does not include the primary publications or a completed prior-art ledger, so neither an exhaustive account of that background nor the novelty of the candidate interaction can yet be claimed.

The central distinction is between failure to produce a regenerative epidermal transition and failure to couple an early injury response to that transition. Metabolic state, tissue mixture, batch structure, and ordinary developmental differences provide competing explanations. A useful candidate must specify a stage-dependent relationship beyond those alternatives, not merely recover the labels used to define competent and refractory groups.

\section{Methods and evidence}
Before analysis, verify the sample manifest and accessible files for GSE245320 and its nominated components: bulk RNA-seq GSE245317/GSE245318, single-cell RNA-seq GSE245312, and spatial data GSE245313. Confirm developmental stages, actual sampling times, animal or pool identities, biological replication, technical replication, batch, tissue boundaries, and any linked regeneration phenotypes. Treat the reported stage-40/stage-46 comparison and sampling at zero, half an hour, six hours, and three days after amputation as starting information to verify, not as a completed inventory.

Define the biological replicate before constructing expression summaries. For temporal coupling, establish whether early and later measurements can be linked through an appropriate biological design; do not pair unrelated destructive samples, replicate numbers, or stage averages as though they were individual trajectories. If the design supports only cross-sectional stage-by-time contrasts, report those contrasts without interpreting them as an animal-level early-to-late predictor.

Predefine early wound-response and later RIC/ECM summaries, then compare the proposed stage-dependent relationship against stage-only, RIC/ECM-only, and known-signaling baselines. Assess batch structure, multicollinearity, replicate sensitivity, multiple testing, and uncertainty at the biological-replicate level. The exact scoring definitions, model specification, performance metric, and numerical decision thresholds remain to be preregistered; no fitted comparison is available in the supplied evidence.

Use single-cell and spatial information, where coverage permits, to distinguish epidermal state changes from changes in cell abundance, tissue mixture, or generalized wound healing. Treat these modalities as complementary observations from a shared study, not independent replication. Plan stage-matched rescue and blocking tests with sham, injury, timing, toxicity, and target-engagement controls, followed by validation on previously unseen biological material; none of these experiments is reported as performed.

\section{The developmental switch and the mechanistic distinction}
The explanatory target is the developmental loss and subsequent recovery of tail-regeneration competence in Xenopus laevis tadpoles, rather than the separation of two developmental-stage labels. Because competence varies between animals, a molecular difference between stage 40 and stage 46 would not, by itself, establish a predictor of an individual animal’s regenerative outcome. A useful mechanistic account must instead identify a conditional relationship that can be tested against regeneration outcomes and explain why that relationship changes across development, including when competence returns.

The biological starting point is the reported transient formation of regeneration-initiating cells (RICs) from basal epidermal cells and the association of extracellular-matrix (ECM) remodeling defects with regeneration failure in the 2024 Genome Biology study (doi:10.1186/s13059-024-03396-3). RIC formation and ECM remodeling therefore constitute prior-art background, not the proposed discovery. Likewise, established associations involving reactive oxygen species, inflammatory signals, and developmental signaling motivate candidate inputs but do not themselves explain the developmental switch.

Nutritional influences reported in Xenopus tropicalis provide an alternative hypothesis to examine, not evidence that the same stage-dependent mechanism operates in X. laevis.

The central distinction is between failure to execute a regenerative epidermal transition and failure to couple an earlier injury response to that transition. Under the first explanation, the later RIC/ECM program could remain deficient even when the early wound response is relatively strong. Under the second, developmental stage changes the relationship between early wound-response intensity and subsequent activation of the regenerative program. Both explanations could produce lower later-program activity at a refractory stage; their distinguishing prediction concerns whether later activation depends on the preceding injury response differently across stages, rather than whether either program differs in average abundance.

The prioritized hypothesis is that early wound-response intensity is positively associated with later epidermal RIC/ECM activation at stage 40, whereas this association is weaker at stage 46. The intended early window extends from half an hour to six hours after amputation, with later activation assessed at the reported three-day sampling point, subject to verification of sample coverage and whether that timing adequately captures the transition. Let $\beta_{40}$ and $\beta_{46}$ denote the slopes relating the same predefined early-response summary to the same later-program summary within their respective stages.

The directional prediction is $\beta_{40}>0$ and $\beta_{46}<\beta_{40}$; these inequalities specify a proposed test, not measured effects.

This formulation makes a narrower claim than a general failure of regeneration signaling. Lower later-program activity without evidence of altered coupling would not specifically support the coupling hypothesis over failure of the epidermal transition itself. Metabolic state or tissue composition could instead condition both summaries, while developmental differences or analytical confounding could generate an apparent stage-dependent relationship. The proposed interaction must therefore be evaluated against stage-only, RIC/ECM-only, and familiar signaling-module explanations, rather than treated as informative merely because it contains more variables.

Temporal ordering also must not be mistaken for an animal-level trajectory. The reported early and late transcriptomic sampling points do not establish that an early measurement and a later measurement can be linked within the same animal or biological replicate. Unless such linkage is verified or obtained prospectively, differences between stage-level time courses cannot establish individual injury-to-regeneration coupling. Biological replicates, rather than cells from a shared animal or library, must support the inference, and independent outcome-linked observations are needed before describing the relationship as an animal-level predictor.

The resulting proposal is thus a bounded mechanistic hypothesis, not a demonstrated new mechanism. Its novelty remains unresolved without a traceable assessment of prior and subsequent literature, and its biological interpretation requires independent validation. Most importantly, a stage-40 versus stage-46 contrast cannot explain later regained competence on its own: the same coupling prediction must also be tested at a later competent stage. Evidence that the proposed relationship returns with competence would support an account of the developmental switch; failure to recover that relationship would limit the hypothesis even if it distinguished the initial two stages.

\section{An estimable test of temporal coupling}
The proposed test concerns a stage-dependent relationship between injury sensing and a subsequent regenerative program, rather than the abundance of either program alone. In Xenopus laevis, the hypothesis is that stronger early wound responses are associated with greater later activation of an epidermal regeneration-initiating cell (RIC) and extracellular-matrix (ECM) program at stage 40, whereas this association is weaker at stage 46. Denoting the corresponding slopes by $\beta_{40}$ and $\beta_{46}$, the directional predictions are $\beta_{40}>0$ and $\beta_{46}<\beta_{40}$. These inequalities specify a hypothesis, not estimated effects or an established mechanism.

The intended early-response window spans half an hour to six hours after amputation, with the reported zero-hour samples providing a baseline where the design permits. Later RIC/ECM expression would be assessed at the reported three-day endpoint. The same definitions of early wound-response intensity and later program activation would be applied across stages, with their construction specified before testing the interaction. This endpoint would test an association with later program expression; it would not, by itself, resolve the timing of the intervening transient epidermal transition.

Whether this relationship is estimable depends first on the actual sampling design. The nominated starting material is GEO series GSE245320, including bulk RNA-sequencing series GSE245317 and GSE245318, single-cell series GSE245312, and spatial series GSE245313. Although the published comparison reports stages 40 and 46 and the proposed sampling times, the supplied evidence does not establish the available sample counts, biological replication, measurement linkage, or regeneration-outcome labels. A sample-level manifest must therefore establish developmental stage, collection time, tissue, animal or pool identity, library identity, and batch before any coupling model is fitted.

The biological replicate, rather than the sequencing observation, must determine the unit of inference. If a library represents one animal, its cells or spatial observations cannot be treated as independent animal replicates; if it represents a pool, the inferential unit is that pool unless individual contributions are separately recoverable. Multiple libraries derived from the same biological material must likewise retain their shared identity. The component series belong to the same study and cannot be counted as independent biological replications merely because they use different technologies.

Most importantly, early and late measurements require a justified biological linkage. Destructive samples collected from different animals at different times cannot be paired retrospectively by replicate number, expression similarity, or developmental stage to manufacture individual trajectories. Such a design may estimate stage-specific mean responses over time, but it cannot identify whether an animal with a stronger early response subsequently develops a stronger RIC/ECM program. A pool-level or other linked-unit relationship would have to be named and interpreted at that level, rather than presented as an individual-animal predictor.

If valid linked biological units are available, a simple interaction model would make the hypothesis explicit: \$L\_i=\textbackslash{}alpha+\textbackslash{}gamma S\_i+\textbackslash{}beta E\_i+\textbackslash{}delta S\_iE\_i+\textbackslash{}varepsilon\_i\$, where $L_i$ is later RIC/ECM expression, $E_i$ is the prespecified early-response summary, and $S_i$ indicates stage 46 rather than stage 40. In this parameterization, $\beta_{40}=\beta$ and \$\textbackslash{}beta\_\{46\}=\textbackslash{}beta+\textbackslash{}delta\$, so the directional test concerns $\beta>0$ and \$\textbackslash{}delta\textless{}0\$. Batch and other justified design terms would be incorporated only where distinguishable from stage, with uncertainty calculated at the biological-replicate level. Estimates and intervals for both slopes and their difference would be more informative than an interaction significance test alone.

If the samples are unlinked, the feasible analysis must instead be restricted to stage-by-time contrasts. These could ask whether the early response differs between stages and whether the later RIC/ECM program shows a corresponding stage difference, using independently replicated samples at each observed time. Even concordant contrasts would remain compatible with several explanations, including failure to induce the later program, developmental differences in tissue composition, and stage-specific injury-to-program coupling. They would motivate a linked-measurement experiment but would not constitute an estimate of the proposed early-to-late slope.

The tissue interpretation requires similar restraint. Because the single-cell evidence predominantly characterizes regenerating animals rather than matched snapshots across both conditions, it may inform the identity of an epidermal program without supplying a direct single-cell test of the stage interaction. Bulk RIC/ECM expression must therefore be examined for sensitivity to tissue or cell-composition differences wherever the available measurements support that distinction. A stage-associated change in a bulk program score should not automatically be interpreted as altered transcription within basal epidermal cells, nor should an association that disappears under supported composition controls be retained as evidence for uncoupling.

Model comparisons must distinguish explanation of later expression from prediction of regeneration. For later-program expression, the interaction should first be compared with stage-only and additive early-response models, and then with models incorporating familiar injury-signaling summaries; a stage difference alone would not establish coupling. If verified regeneration outcomes are available, a separate predictive analysis should compare stage-only, RIC/ECM-only, and known-signaling baselines before testing whether coupling adds information on held-out biological units.

Resampling must preserve animal or pool identity and relevant batch structure, while alternative windows and multiple module tests require explicit control of analytic multiplicity.

A reproducible reduction in the early-to-late slope at stage 46 would support the proposed conditional association more specifically than a lower late-program mean alone, but it would not establish a causal failure of signal transmission. The interpretation would weaken if the interaction were unstable across biological replicates, inseparable from batch or composition, or unable to improve on simple baselines. Conversely, failure to obtain a precise positive slope at stage 40 would leave the directional hypothesis unsupported rather than demonstrate that coupling is absent.

No incremental predictive value or mechanistic confirmation has yet been demonstrated; the present contribution is an estimable hypothesis with a clear boundary between a test of coupling and a comparison of stage-associated expression patterns.

\section{Separating coupling failure from cellular and tissue-level alternatives}
The proposed coupling mechanism requires a distinction between activation within the wound epidermis and changes in the amount or composition of tissue sampled. In Xenopus laevis, the working hypothesis predicts that early wound-response intensity is positively associated with subsequent epidermal regeneration-initiating cell (RIC) and extracellular-matrix (ECM) program activation at stage $40$, but less strongly associated at stage $46$. Expressed as stage-specific slopes relating the same early-response summary to the same later-program summary, the prediction is $\beta_{40}>0$ and $\beta_{46}<\beta_{40}$.

These are prospective inequalities, not estimated effects; the intended comparison spans the reported early samples at $0.5$ and $6$ hours after amputation and the later sample at $3$ days, subject to verification of actual coverage.

The spatial dataset GSE245313 and single-cell dataset GSE245312 would be used to determine what biological quantity a candidate bulk signature represents. The analysis would distinguish increased expression of the later program within an epidermal population from increased abundance of RIC-like cells, altered representation of other tissues, or a broader healing response. Marker panels would be fixed before testing the stage-dependent association, with separate summaries for epidermal identity, the candidate transition, and other represented tissue compartments.

The relevant spatial prediction is localization of the candidate later program to the wound-associated epidermal compartment, rather than a tissue-wide expression change interpreted as an epidermal transition.

Single-cell reference information would help identify plausible cellular contributors, but it cannot supply a matched refractory comparison that the available sampling does not provide. GSE245312 predominantly characterizes regenerating animals, so a state recognized there could be sought in the spatial and bulk evidence without treating its apparent absence elsewhere as proof that refractory epidermal cells cannot enter that state. Any reference-based abundance estimate would therefore be treated as conditional on the states represented in the reference. Because these modalities originate from the same study, agreement among them would strengthen localization and interpretation, not constitute independent biological replication.

The composition test would compare the prospective coupling association before and after adjustment for estimable tissue and cell-population contributions. In particular, the analysis would ask whether the stage-dependent association remains when epidermal representation and RIC-like abundance are considered separately from program expression within those compartments. Biological specimens, rather than individual cells or spatial locations, would remain the units of inference; aggregation and uncertainty estimates would follow the verified specimen and library structure.

Cross-time associations would be estimated only at the level supported by the sampling design, and unlinked early and late samples would not be represented as longitudinal measurements from individual tadpoles.

If the prospective association disappears after accounting for tissue mixture, its interpretation as a distinct injury-to-transcription coupling mechanism would be weakened. A signature explained entirely by the amount of epidermis or the abundance of RIC-like cells would not establish altered transcriptional responsiveness within epidermal cells. Nevertheless, abundance adjustment is not by itself a decisive causal test: if injury signaling changes entry into a transient state, the abundance of that state could lie downstream of coupling rather than merely confound its measurement.

The resulting interpretation would depend on whether spatially localized state activation can be separated from population abundance; if it cannot, the evidence would support a tissue-level association rather than a resolved cell-state mechanism.

Generalized healing and developmental or metabolic differences would remain explicit alternatives even if the signal localizes to the epidermis. The candidate association would therefore be evaluated against simpler stage-only, RIC/ECM-only, and known injury-signaling summaries, with attention to whether it adds information beyond those baselines. Where suitable measures are available, the analysis would ask whether a general healing signature accounts for the same pattern without requiring a stage-specific coupling rule. Metabolic gating would remain unresolved if the available data do not measure it adequately; observations from X. tropicalis would not be treated as evidence that the same explanation applies to X. laevis.

A proposed causal test would distinguish restoring the amplitude of an early injury signal, restoring its transmission to the epidermal program, and directly enabling the later transition. These interventions are not interchangeable. If stage $46$ animals have an insufficient early input but retain a functional response pathway, restoring that input should restore later activation; if transmission is impaired, increasing the input alone need not do so. A rescue after direct induction of the later transition would demonstrate bypass under the tested conditions, but would not by itself distinguish defective coupling from failure to initiate that transition.

The more discriminating rescue-versus-block experiment would require a selective means of restoring the proposed injury-to-epidermal transmission step at stage $46$, together with a corresponding block at stage $40$. Under a coupling-failure model, the rescue should restore the later epidermal program and improve regrowth, whereas the block should suppress the later program despite a verified early injury response. Under an intrinsic-transition-failure model, restoring upstream transmission should fail to restore the later program even when target engagement is demonstrated.

A separate direct-transition arm would test downstream responsiveness, but a failed rescue would remain uninterpretable if the intervention had not demonstrably acted on its intended step.

Each perturbation would require temporally resolved measurements of the early response, epidermal transition, ECM remodeling, and subsequent regrowth, rather than regrowth alone. Stage-matched sham, injury-only, and vehicle or procedure controls would separate treatment effects from injury and handling, while toxicity and developmental progression would be assessed alongside regeneration. The primary regrowth endpoint would be a blinded quantitative morphological measure at a prespecified post-amputation time, with randomization and sample-size assumptions defined at the animal level.

These are proposed design requirements: the supplied evidence identifies no selective target, validated reagent, dose, or intervention window, so it does not yet support an executable decisive perturbation protocol.

The bounded outcome of the observational analysis is therefore a test of whether the proposed stage-dependent association survives localization and composition alternatives, not a demonstration of causal coupling. Loss of the association after those tests would argue against advancing it as a distinct within-epidermal predictor; persistence would justify developing the selective perturbations needed to separate transmission failure from transition failure. Independent biological evidence from new specimens or a validated perturbation would be required before interpreting the proposed rule as a new regeneration mechanism.

\section{Validation, later recovery, and the limits of the present claim}
Validation must distinguish a previously unreported hypothesis, a reproducible computational association, and a biological mechanism. For the proposed injury-to-regeneration coupling rule, the relevant computational claim would be that early wound-response intensity carries information about subsequent epidermal regeneration-initiating cell (RIC) and extracellular-matrix (ECM) activation beyond developmental stage and established regenerative programs. A mechanistic claim would additionally require independent biological evidence that discriminates a coupling defect from failure of the downstream transition itself. These evidentiary levels should remain separate even if the same measurements contribute to each.

The novelty assessment should therefore address the conditional relationship, not merely its component pathways. The supplied background attributes basal-epidermal RIC origins and ECM-remodeling defects associated with regeneration failure to the 2024 Genome Biology study identified by DOI 10.1186/s13059-024-03396-3; those observations constitute prior art rather than the proposed discovery. A primary-source ledger, updated through the investigation date, should record the organism, developmental stages, injury timing, tissue, measured outcome, and associative or perturbational basis of each relevant claim.

It should explicitly determine whether a stage-dependent relationship between early injury responses and later RIC/ECM activation has already been tested, while preserving conflicting findings and differences in experimental context. Observations concerning nutritional state in Xenopus tropicalis should be recorded separately and should not establish either stage behavior or a metabolic mechanism in Xenopus laevis.

Computational validation should begin with prespecified stage-only, RIC/ECM-only, and familiar injury-signaling baselines. The coupling model should be evaluated by the numerical direction and magnitude of its proposed stage-dependent association, its uncertainty, and its incremental performance on held-out biological units, rather than by an apparent improvement in the discovery samples alone. Biological replicates must remain the units of inference and resampling; cells or spatial measurements from the same animal or library cannot supply independent replication.

Sensitivity analyses should assess whether the relationship persists under batch-aware resampling, alternative defensible time windows, and adjustment for tissue composition, with multiplicity controlled where several candidate relationships are tested. An unstable estimate or a result explained by composition or batch would not support the proposed interaction.

The endpoint determines what such validation could establish. Distinguishing stage $40$ from stage $46$ is not equivalent to predicting whether an individual tadpole regenerates, particularly because competence varies between animals. Animal-level prediction would require measured regeneration outcomes linked to the relevant biological units; absent those labels, an analysis should remain a stage-associated transcriptional comparison. Likewise, early and late samples should not be paired analytically as though they came from the same animal unless their provenance supports that relationship. A temporal pattern across sampled groups can motivate a coupling hypothesis without demonstrating an individual-level dependency.

Independent validation should use previously unseen X. laevis biological material or an independently generated perturbation-and-phenotype dataset, with the candidate relationship and evaluation criteria fixed before examining validation outcomes. The bulk, single-cell, and spatial components of GSE245320 share a study and therefore cannot be counted as independent biological replications merely because they use different assays. Their complementary measurements may help assess tissue localization or composition alternatives, but confirmation of the proposed prediction must extend beyond reuse of the discovery material.

To support a mechanism rather than an association, rescue-versus-block evidence would need to distinguish restoration or disruption of injury-to-program coupling from a nonspecific change in downstream RIC/ECM activation, using appropriate controls and measured regeneration outcomes.

The later return of competence provides a further explanatory test. A comparison restricted to stages $40$ and $46$ cannot establish why regeneration subsequently becomes possible again. If disrupted coupling explains the temporary loss, the proposed prediction is that a later, independently verified competent X. laevis stage should restore the relationship between early wound responses and subsequent epidermal RIC/ECM activation. This test should retain the prespecified temporal ordering and tissue interpretation while comparing the coupling rule with the same simpler baselines.

Later regeneration in the absence of restored coupling would limit the rule as an explanation of the full developmental switch, even if it remained associated with the earlier contrast.

The present evidence supplies no completed primary-source novelty assessment, reproducible computational association, numerical effect with uncertainty, or independent biological confirmation. This absence is not an experimental null result and does not reject any of the competing biological hypotheses. The defensible conclusion is therefore a hypothesis-led validation plan: a computational result would require reproducible improvement over the prespecified baselines on held-out biological units, whereas a new mechanistic explanation would require discriminating biological confirmation and an account of later regained competence.

Until those conditions are met, stage-dependent injury-to-regeneration coupling remains an unvalidated candidate whose novelty and explanatory scope are unresolved.

\section{Limitations}
\begin{itemize}
\item No completed biological or computational result is available in the supplied selection; numerical effects, uncertainty, predictive performance, and intervention outcomes cannot be reported.
\item The candidate coupling interaction originated in the project hypothesis space. Its presence in a research strategy does not establish novelty or empirical support.
\item Accession-level coverage, biological replication, sample linkage, batch structure, and animal-level regeneration labels have not been verified.
\item The nominated single-cell and spatial resources share study provenance, and the described single-cell coverage does not provide a matched refractory atlas.
\item A two-stage comparison cannot explain later recovery without additional observations, and Xenopus tropicalis findings cannot substitute for within-species validation in Xenopus laevis.
\item The bounded evidence index is not an exhaustive review of saved material or the literature; its omitted notes cannot be treated as inspected evidence.
\end{itemize}
\section{Open obligations}
\begin{itemize}
\item Retrieve and verify primary publications, including subsequent work through the investigation date, and build a cited ledger that explicitly searches for the proposed stage-dependent temporal interaction.
\item Inspect GEO records and files to establish actual stages, times, biological units, replicate counts, batches, tissue sampling, data availability, and phenotype linkage.
\item Determine whether temporal coupling is identifiable from the available sampling design; otherwise restrict the retrospective analysis to justified cross-sectional contrasts and design a prospective linkage study.
\item Preregister gene-module definitions, timing contrasts, model terms, uncertainty estimation, multiple-testing procedures, negative controls, performance metrics, and stopping rules.
\item Estimate a numerical coupling difference and uncertainty, compare it with all specified simple baselines, and test sensitivity to biological replicates, batch, composition, and alternate time windows.
\item Identify a selective, feasible perturbation and specify its target engagement, timing, counterfactual predictions, sham and injury controls, toxicity monitoring, blinded endpoint, sample-size assumptions, and animal-use requirements.
\item Validate the locked prediction on previously unseen specimens or independently generated conditions without using validation outcomes to select the signature.
\item Test a later competent Xenopus laevis stage to determine whether recovery accompanies restoration of the proposed coupling.
\end{itemize}
\section{Conclusions}
Stage-dependent injury-to-regeneration coupling is a specific candidate explanation to test, not an established explanation of the refractory period. Its proposed contribution is an interaction between early injury response and later epidermal activation, rather than rediscovery of RICs or ECM remodeling.

The first decisive issue is identifiability: verified sample structure must determine whether an early-to-late relationship can be estimated at all. A stage-associated expression pattern must not be promoted to an animal-level regeneration predictor.

The available evidence supports neither a new mechanism nor a completed computational result. A defensible advance requires a verified novelty boundary, quantitative improvement over baselines, discrimination from alternative explanations, and independent biological support that also addresses later recovery.

\section*{Authorship note}
AI-assisted synthesis from saved research; not independently peer reviewed.
\end{document}
